% fixture: amsart/amsart-theorems
% tier: T3
% engine: document
% covers: \vspace* at the start (\prevdepth is -1000pt there in amsart: no \onecolumn), amsart's built-in amsthm: plain, definition and remark styles, numbering within the section, a note, the proof with amsart's \qedsymbol (\textsquare from msam)
% issue: #360
\documentclass{amsart}
\newtheorem{theorem}{Theorem}[section]
\newtheorem{lemma}[theorem]{Lemma}
\theoremstyle{definition}
\newtheorem{definition}[theorem]{Definition}
\theoremstyle{remark}
\newtheorem*{remark}{Remark}
\begin{document}
\vspace*{1cm}
\section{Results}
Some words before the first theorem of the section.
\begin{theorem}[Fermat]
Every prime congruent to one modulo four is a sum of two squares, and the representation is unique up to order and signs.
\end{theorem}
\begin{proof}
Consider the involution on the finite set of solutions and count its fixed points carefully.
\end{proof}
\begin{definition}
A number is \emph{representable} if it is a sum of two squares of integers.
\end{definition}
\begin{lemma}
The product of two representable numbers is representable.
\end{lemma}
\begin{remark}
The identity behind the lemma is due to Brahmagupta and Fibonacci.
\end{remark}
Text after the remark.
\end{document}
